% ============================================================
%  第 1 章  坐标系、矢量分析与张量初步
%  教材完整版：~50页
% ============================================================

\chapter{坐标系、矢量分析与张量初步}

\section{本章概览与学习路线}

\subsection{本章在物理中的位置}

物理定律的数学表述离不开"空间"和"变化"两个概念。牛顿第二定律 $\mathbf{F}=m\mathbf{a}$ 中的加速度 $\mathbf{a}$ 是在某个坐标系中定义的二阶时间导数；麦克斯韦方程组是矢量场 $\mathbf{E}$ 和 $\mathbf{B}$ 在三维空间中的微积分关系；薛定谔方程中的哈密顿量 $\hat{H}$ 是线性空间上的厄米算符。所有这些物理理论，共享同一套数学语言——矢量分析、张量、线性空间。

本章构建这套语言。你将学到：

\begin{itemize}
  \item 如何在任意曲线坐标系中写出物理定律（拉梅系数）
  \item 矢量微积分的三大核心算符（梯度、散度、旋度）
  \item 连接局域与整体的积分定理（高斯、斯托克斯、格林）
  \item 矢量场分解的根本定理（亥姆霍兹分解）
  \item 保证物理定律坐标无关性的工具（张量）
  \item 量子力学的数学舞台（线性空间与厄米算符）
\end{itemize}

\subsection{前置知识}

阅读本章前，你应该熟悉：
\begin{itemize}
  \item 一元和多元微积分（偏导数、重积分、线面积分）
  \item 线性代数基础（矩阵乘法、行列式、本征值）
  \item 普通物理中的矢量运算（点积、叉积）
\end{itemize}

\subsection{学习路径建议}

\begin{itemize}
  \item \textbf{物理专业学生}：重点在 \S1.3--\S1.6（矢量分析）和 \S1.8--\S1.9（线性空间），这是电动力学和量子力学的直接预备。
  \item \textbf{数学专业学生}：重点在 \S1.7（张量）和 \S1.9（谱定理），这些是微分几何和泛函分析的入门。
  \item \textbf{工程专业学生}：重点在 \S1.3--\S1.6 和 \S1.7（转动惯量、应力张量），为连续介质力学打基础。
\end{itemize}

\section{引言：物理定律需要坐标系吗？}

\subsection{一个基本问题}

牛顿的推理是这样的：如果一个物体不受外力，"它将继续保持静止或匀速直线运动"。但问题在于——\textit{相对于谁静止？}

你在匀速行驶的火车上捡起一枚硬币，它做竖直直线运动（相对于火车），同时做抛物线运动（相对于地面）。两个描述都是正确的，但用的坐标系不同。这个看似平凡的事实引出了一个深刻的物理原理：\textbf{物理定律的形式在所有惯性系中相同}（伽利略相对性原理）。

但数学上，"同一物理定律在不同坐标系下形式不变"是什么意思？以梯度为例：
\[
\nabla f = \frac{\partial f}{\partial x}\mathbf{e}_x + \frac{\partial f}{\partial y}\mathbf{e}_y + \frac{\partial f}{\partial z}\mathbf{e}_z
\]
在直角坐标中是这样的形式。但在球坐标中，我们将会看到
\[
\nabla f = \frac{\partial f}{\partial r}\mathbf{e}_r + \frac{1}{r}\frac{\partial f}{\partial\theta}\mathbf{e}_\theta + \frac{1}{r\sin\theta}\frac{\partial f}{\partial\phi}\mathbf{e}_\phi .
\]

看起来不同，但\textbf{物理量是同一个 $\nabla f$}，只是表达方式不同。理解这种"不同写法代表同一物理量"的能力，是本章的第一个核心目标。

\subsection{本章的逻辑主线}

\begin{center}
\fbox{\parbox{0.85\textwidth}{
\textbf{物理问题 $\to$ 选择坐标系 $\to$ 拉梅系数描述几何 $\to$ 微分算符表达定律 $\to$ 积分定理揭示守恒 $\to$ 张量保证协变 $\to$ 线性空间抽象 $\to$ 本征值连接可观测量}
}}
\end{center}

这条线索将在后续所有章节中被反复强化。每当你在第 2 章遇到复变积分、在第 3 章遇到傅里叶变换、在第 4 章遇到 PDE，你都会回到本章建立的框架——\textbf{物理对称性决定数学形式，数学形式揭示物理规律}。

\section{正交曲线坐标系与拉梅系数}

\subsection{为什么不能只用直角坐标？}

直角坐标 $(x,y,z)$ 是可以描述任何点的，但在物理学中，我们遇到的绝大多数问题都具有某种对称性：

\begin{itemize}
  \item \textbf{球对称}：中心力场、氢原子、点电荷的电场
  \item \textbf{柱对称}：无限长导线、圆波导、光纤
  \item \textbf{抛物线对称}：匀强电场中的带电粒子轨迹
\end{itemize}

如果将一个球对称问题强制塞进直角坐标，边界条件会变得极其复杂。例如球的表面 $x^2+y^2+z^2=R^2$ 需要同时约束三个坐标。而在球坐标中，边界只是简洁的 $r=R$。

\textbf{基本原则：让坐标面与物理边界重合。}

\subsection{从直角坐标到一般曲线坐标}

设三维欧氏空间中有一组新的坐标 $(q_1,q_2,q_3)$，它们与直角坐标 $(x,y,z)$ 之间的变换关系为
\begin{equation}
x = x(q_1,q_2,q_3),\quad
y = y(q_1,q_2,q_3),\quad
z = z(q_1,q_2,q_3),
\label{eq:general_transform}
\end{equation}
其中变换函数是光滑的，且雅可比矩阵 $\partial(x,y,z)/\partial(q_1,q_2,q_3)$ 在讨论区域内非奇异（保证变换可逆）。

\subsection{拉梅系数的几何意义}

考虑一个无限小的位移：
\[
\dd\mathbf{r} = \frac{\partial\mathbf{r}}{\partial q_1}\dd q_1
+ \frac{\partial\mathbf{r}}{\partial q_2}\dd q_2
+ \frac{\partial\mathbf{r}}{\partial q_3}\dd q_3 .
\]

每个 $\partial\mathbf{r}/\partial q_i$ 是一个矢量，指向 $q_i$ 增加的方向。其长度
\[
h_i = \left|\frac{\partial\mathbf{r}}{\partial q_i}\right|,\qquad i=1,2,3
\]
称为\textbf{拉梅系数 (Lamé coefficients)}。沿 $q_i$ 方向的单位矢量为
\[
\mathbf{e}_i = \frac{1}{h_i}\frac{\partial\mathbf{r}}{\partial q_i}.
\]

\textbf{弧长平方：}
\[
\dd s^2 = \dd\mathbf{r}\cdot\dd\mathbf{r}
= \sum_{i,j} \frac{\partial\mathbf{r}}{\partial q_i}\cdot\frac{\partial\mathbf{r}}{\partial q_j}\dd q_i\dd q_j .
\]

如果 $\mathbf{e}_i\cdot\mathbf{e}_j = \delta_{ij}$（即坐标线处处正交），则交叉项（$i\neq j$ 的项）消失：
\[
\boxed{\;\dd s^2 = h_1^2\dd q_1^2 + h_2^2\dd q_2^2 + h_3^2\dd q_3^2\;}.
\label{eq:metric_line_element}
\]

这就是\textbf{正交曲线坐标系}的定义。

\begin{example}[完整推导：三种坐标系的拉梅系数]
\textbf{(1) 柱坐标 $(\rho,\phi,z)$}
\[
x = \rho\cos\phi,\quad y = \rho\sin\phi,\quad z = z.
\]

计算偏导：
\[
\frac{\partial\mathbf{r}}{\partial\rho} = (\cos\phi,\sin\phi,0),\quad
\frac{\partial\mathbf{r}}{\partial\phi} = (-\rho\sin\phi,\rho\cos\phi,0),\quad
\frac{\partial\mathbf{r}}{\partial z} = (0,0,1).
\]

取模：
\[
h_\rho = 1,\quad h_\phi = \sqrt{\rho^2\sin^2\phi+\rho^2\cos^2\phi} = \rho,\quad h_z = 1.
\]

弧长元：
\[
\dd s^2 = \dd\rho^2 + \rho^2\dd\phi^2 + \dd z^2.
\]

\textbf{(2) 球坐标 $(r,\theta,\phi)$}
\[
x = r\sin\theta\cos\phi,\quad y = r\sin\theta\sin\phi,\quad z = r\cos\theta.
\]

偏导数：
\begin{align*}
\frac{\partial\mathbf{r}}{\partial r} &= (\sin\theta\cos\phi,\sin\theta\sin\phi,\cos\theta),\\
\frac{\partial\mathbf{r}}{\partial\theta} &= (r\cos\theta\cos\phi,r\cos\theta\sin\phi,-r\sin\theta),\\
\frac{\partial\mathbf{r}}{\partial\phi} &= (-r\sin\theta\sin\phi,r\sin\theta\cos\phi,0).
\end{align*}

取模：
\[
h_r = 1,\quad h_\theta = r,\quad h_\phi = r\sin\theta.
\]

弧长元：
\[
\dd s^2 = \dd r^2 + r^2\dd\theta^2 + r^2\sin^2\theta\,\dd\phi^2.
\]

\textbf{(3) 抛物线坐标 $(\xi,\eta,\phi)$}
\[
x = \xi\eta\cos\phi,\quad y = \xi\eta\sin\phi,\quad z = \frac12(\xi^2-\eta^2).
\]

拉梅系数：
\[
h_\xi = h_\eta = \sqrt{\xi^2+\eta^2},\quad h_\phi = \xi\eta.
\]

这组坐标在带电粒子在匀强电场中的运动问题中极为有用。
\end{example}

\begin{insight}[拉梅系数 = 局部标尺]
拉梅系数的物理图像非常直观：当你沿 $q_i$ 方向走一个坐标步 $\dd q_i$ 时，实际在三维空间中移动了 $h_i\dd q_i$ 这么远。因此 $h_i$ 就是"坐标步长到实际距离的换算因子"。在球坐标中，沿 $\theta$ 方向走 $\dd\theta$ 对应实际距离 $r\dd\theta$（$h_\theta=r$），沿 $\phi$ 方向走 $\dd\phi$ 对应实际距离 $r\sin\theta\,\dd\phi$（$h_\phi=r\sin\theta$）。你可以把这个想象成地球上的经纬度——经度一度对应的地面距离取决于你所在的纬度。
\end{insight}

\section{微分算符在曲线坐标系中的形式}

\subsection{梯度}

梯度 $\nabla f$ 的物理定义是"方向导数最大的方向"及其"最大变化率"。在曲线坐标中，沿 $q_i$ 方向的方向导数为 $(1/h_i)\partial f/\partial q_i$，因此
\[
\boxed{\;\nabla f = \sum_{i=1}^3 \frac{1}{h_i}\frac{\partial f}{\partial q_i}\mathbf{e}_i\;}.
\label{eq:grad_curvilinear}
\]

\subsection{散度}

散度 $\nabla\cdot\mathbf{A}$ 的物理定义是"单位体积的通量"。考虑一个由坐标线围成的无限小曲边六面体，穿过两个相对面的净通量差给出：
\[
\boxed{\;
\nabla\cdot\mathbf{A}
= \frac{1}{h_1h_2h_3}\sum_{i=1}^3
\frac{\partial}{\partial q_i}\left(\frac{h_1h_2h_3}{h_i}A_i\right)
\;}.
\label{eq:div_curvilinear}
\]

\begin{example}[柱坐标中的散度]
$h_\rho=1,\;h_\phi=\rho,\;h_z=1$，$h_1h_2h_3 = \rho$。
\begin{align*}
\nabla\cdot\mathbf{A}
&= \frac{1}{\rho}\left[
\frac{\partial}{\partial\rho}(\rho A_\rho)
+ \frac{\partial}{\partial\phi}(A_\phi)
+ \frac{\partial}{\partial z}(\rho A_z)
\right] \\
&= \frac{1}{\rho}\frac{\partial}{\partial\rho}(\rho A_\rho)
+ \frac{1}{\rho}\frac{\partial A_\phi}{\partial\phi}
+ \frac{\partial A_z}{\partial z}.
\end{align*}
\end{example}

\subsection{旋度}

旋度 $\nabla\times\mathbf{A}$ 的物理定义是"单位面积的最大环量"。同样从六面体的面环量出发：
\[
\boxed{\;
\nabla\times\mathbf{A}
= \sum_{i=1}^3\frac{1}{h_jh_k}
\left[\frac{\partial}{\partial q_j}(h_kA_k)-\frac{\partial}{\partial q_k}(h_jA_j)\right]\mathbf{e}_i
\;}.
\label{eq:curl_curvilinear}
\]

其中 $(i,j,k)$ 取 $(1,2,3)$ 的循环置换。

\begin{example}[球坐标中的旋度]
\begin{align*}
\nabla\times\mathbf{A}
&= \frac{1}{r\sin\theta}\left[
\frac{\partial}{\partial\theta}(\sin\theta A_\phi)
- \frac{\partial A_\theta}{\partial\phi}
\right]\mathbf{e}_r \\
&\quad + \frac{1}{r}\left[
\frac{1}{\sin\theta}\frac{\partial A_r}{\partial\phi}
- \frac{\partial}{\partial r}(r A_\phi)
\right]\mathbf{e}_\theta \\
&\quad + \frac{1}{r}\left[
\frac{\partial}{\partial r}(r A_\theta)
- \frac{\partial A_r}{\partial\theta}
\right]\mathbf{e}_\phi .
\end{align*}
\end{example}

\subsection{拉普拉斯算子}

拉普拉斯算子 $\nabla^2 f = \nabla\cdot(\nabla f)$，将梯度公式代入散度公式：
\[
\boxed{\;
\nabla^2 f = \frac{1}{h_1h_2h_3}\sum_{i=1}^3
\frac{\partial}{\partial q_i}\left(\frac{h_1h_2h_3}{h_i^2}
\frac{\partial f}{\partial q_i}\right)
\;}.
\label{eq:laplacian_curvilinear}
\]

\begin{example}[完整算例：氢原子的拉普拉斯算子]
氢原子的薛定谔方程在球坐标中写为：
\[
-\frac{\hbar^2}{2m}\nabla^2\psi - \frac{e^2}{4\pi\varepsilon_0 r}\psi = E\psi.
\]

代入 $h_r=1,\;h_\theta=r,\;h_\phi=r\sin\theta$：
\begin{align*}
\nabla^2\psi &= \frac{1}{r^2}\frac{\partial}{\partial r}\left(r^2\frac{\partial\psi}{\partial r}\right)
+ \frac{1}{r^2\sin\theta}\frac{\partial}{\partial\theta}\left(\sin\theta\frac{\partial\psi}{\partial\theta}\right)
+ \frac{1}{r^2\sin^2\theta}\frac{\partial^2\psi}{\partial\phi^2} \\
&= \frac{1}{r^2}\frac{\partial}{\partial r}\left(r^2\frac{\partial\psi}{\partial r}\right)
- \frac{\hat{L}^2}{\hbar^2 r^2}\psi,
\end{align*}
其中 $\hat{L}^2 = -\hbar^2\left[\frac{1}{\sin\theta}\frac{\partial}{\partial\theta}\left(\sin\theta\frac{\partial}{\partial\theta}\right) + \frac{1}{\sin^2\theta}\frac{\partial^2}{\partial\phi^2}\right]$ 是角动量平方算符。

这正是分离变量法的出发点——径向部分和角向部分自然分离，因为拉普拉斯算子的球坐标形式就包含了这种结构。
\end{example}

\begin{example}[完整算例：点电荷电场的散度验证]
点电荷 $q$ 在球坐标中的电场 $\mathbf{E} = \frac{q}{4\pi\varepsilon_0 r^2}\mathbf{e}_r$。

用散度公式验证 $\nabla\cdot\mathbf{E}=0$（$r>0$）：
\begin{align*}
\nabla\cdot\mathbf{E}
&= \frac{1}{r^2}\frac{\partial}{\partial r}\left(r^2 E_r\right) \\
&= \frac{1}{r^2}\frac{\partial}{\partial r}\left(r^2\cdot\frac{q}{4\pi\varepsilon_0 r^2}\right) \\
&= \frac{1}{r^2}\frac{\partial}{\partial r}\left(\frac{q}{4\pi\varepsilon_0}\right) = 0.
\end{align*}
关键点：$r^2 E_r = q/(4\pi\varepsilon_0)$ 是常数——这正是因为 $\mathbf{E}\propto 1/r^2$。如果电场按 $1/r^{2+\varepsilon}$ 衰减，散度就不为零了。这说明\textbf{库仑平方反比律直接等价于 $\nabla\cdot\mathbf{E}=0$（除点电荷外）}。
\end{example}

\section{Levi-Civita 符号与矢量代数}

\subsection{为什么需要一种新的符号？}

矢量分析中有大量恒等式，用分量展开证明非常冗长。例如
\[
\nabla\times(\nabla\times\mathbf{A}) = \nabla(\nabla\cdot\mathbf{A}) - \nabla^2\mathbf{A},
\]
在电磁学中几乎每天都要用到。如果用分量展开硬算，三个方向共 9 项。而用指标符号只需 3 行。

更重要的是，指标符号是张量分析的入门钥匙——后续的狭义相对论和广义相对论将大量使用这种语言。

\subsection{定义}

\textbf{Levi-Civita 符号} $\epsilon_{ijk}$（也称为排列符号或反对称符号）定义为：
\[
\epsilon_{ijk} =
\begin{cases}
+1, & (i,j,k) \text{ 为 } (1,2,3)\text{ 的偶排列},\\
-1, & (i,j,k) \text{ 为 }(1,2,3)\text{ 的奇排列},\\
0, & \text{有重复指标}.
\end{cases}
\]

偶排列：$(1,2,3), (2,3,1), (3,1,2)$。
奇排列：$(3,2,1), (1,3,2), (2,1,3)$。

\subsection{用 \texorpdfstring{$\epsilon_{ijk}$}{εijk} 表示叉积和行列式}

\textbf{叉积：}
\[
(\mathbf{A}\times\mathbf{B})_i = \epsilon_{ijk}A_jB_k,
\]
其中重复指标 $j,k$ 自动求和（\textbf{爱因斯坦求和约定}）。

\textbf{行列式：}
\[
\det\begin{pmatrix}
a_{11} & a_{12} & a_{13}\\
a_{21} & a_{22} & a_{23}\\
a_{31} & a_{32} & a_{33}
\end{pmatrix}
= \epsilon_{ijk}a_{1i}a_{2j}a_{3k}.
\]

\subsection{最重要的恒等式：\texorpdfstring{$\epsilon$-$\delta$}{ε-δ} 关系}

\[
\boxed{\;\epsilon_{ijk}\epsilon_{ilm} = \delta_{jl}\delta_{km} - \delta_{jm}\delta_{kl}\;}.
\label{eq:epsilon_delta}
\]

这里对 $i$ 求和。理解这个公式：
\begin{itemize}
\item 左边：两个 Levi-Civita 符号的乘积，对共同指标 $i$ 求和。
\item 右边：两个 Kronecker $\delta$ 的乘积之差。
\item 口诀：\textbf{"积的差"}——$\delta_{jl}\delta_{km}$ 减交换 $l,m$ 的 $\delta_{jm}\delta_{kl}$。
\end{itemize}

\subsection{恒等式的指标法证明}

\begin{example}[三重积 $\mathbf{A}\times(\mathbf{B}\times\mathbf{C})$]
\textbf{目标：} $\mathbf{A}\times(\mathbf{B}\times\mathbf{C}) = \mathbf{B}(\mathbf{A}\cdot\mathbf{C}) - \mathbf{C}(\mathbf{A}\cdot\mathbf{B})$。

第 $i$ 分量：
\begin{align*}
[\mathbf{A}\times(\mathbf{B}\times\mathbf{C})]_i
&= \epsilon_{ijk}A_j(\mathbf{B}\times\mathbf{C})_k \\
&= \epsilon_{ijk}A_j\epsilon_{klm}B_lC_m \\
&= \epsilon_{kij}\epsilon_{klm}A_jB_lC_m \\
&= (\delta_{il}\delta_{jm} - \delta_{im}\delta_{jl})A_jB_lC_m.
\end{align*}
展开两项：
\begin{itemize}
\item $\delta_{il}\delta_{jm}A_jB_lC_m = A_jB_iC_j = B_i(\mathbf{A}\cdot\mathbf{C})$。
\item $\delta_{im}\delta_{jl}A_jB_lC_m = A_jB_jC_i = C_i(\mathbf{A}\cdot\mathbf{B})$。
\end{itemize}
所以 $[\mathbf{A}\times(\mathbf{B}\times\mathbf{C})]_i = B_i(\mathbf{A}\cdot\mathbf{C}) - C_i(\mathbf{A}\cdot\mathbf{B})$，即矢量形式。$\square$
\end{example}

\begin{example}[三重积 $(\mathbf{A}\times\mathbf{B})\times\mathbf{C}$]
\[
(\mathbf{A}\times\mathbf{B})\times\mathbf{C} = \mathbf{B}(\mathbf{A}\cdot\mathbf{C}) - \mathbf{A}(\mathbf{B}\cdot\mathbf{C}).
\]
第 $i$ 分量：
\begin{align*}
[(\mathbf{A}\times\mathbf{B})\times\mathbf{C}]_i
&= \epsilon_{ijk}(\mathbf{A}\times\mathbf{B})_jC_k \\
&= \epsilon_{ijk}\epsilon_{jlm}A_lB_mC_k \\
&= -\epsilon_{jik}\epsilon_{jlm}A_lB_mC_k \\
&= -(\delta_{il}\delta_{km} - \delta_{im}\delta_{kl})A_lB_mC_k \\
&= -A_iB_kC_k + A_kB_iC_k \\
&= B_i(\mathbf{A}\cdot\mathbf{C}) - A_i(\mathbf{B}\cdot\mathbf{C}).
\end{align*}
注意结果中 $\mathbf{A}$ 和 $\mathbf{C}$ 的位置与第一个三重积不同——叉积不满足结合律！$\square$
\end{example}

\begin{example}[证明 $\nabla\times(\nabla f) = 0$]
\begin{align*}
[\nabla\times(\nabla f)]_i
&= \epsilon_{ijk}\partial_j\partial_k f.
\end{align*}
由于 $\partial_j\partial_k f = \partial_k\partial_j f$（混合偏导与次序无关），而 $\epsilon_{ijk}$ 关于 $j,k$ 反对称，求和为零。$\square$
\end{example}

\begin{example}[证明 $\nabla\cdot(\nabla\times\mathbf{A}) = 0$]
\begin{align*}
\nabla\cdot(\nabla\times\mathbf{A})
&= \partial_i(\epsilon_{ijk}\partial_j A_k) \\
&= \epsilon_{ijk}\partial_i\partial_j A_k.
\end{align*}
同样 $\partial_i\partial_j = \partial_j\partial_i$ 关于 $i,j$ 对称，$\epsilon_{ijk}$ 关于 $i,j$ 反对称，结果为零。$\square$
\end{example}

\subsection{物理洞察：为什么 \texorpdfstring{$\epsilon_{ijk}$}{εijk} 这么有用？}

物理定律经常涉及叉积。力矩 $\bm{\tau} = \mathbf{r}\times\mathbf{F}$，洛伦兹力 $\mathbf{F}=q\mathbf{v}\times\mathbf{B}$，角动量 $\mathbf{L} = \mathbf{r}\times\mathbf{p}$。用指标符号，这些可以统一处理。

更重要的是，$\epsilon_{ijk}$ 是三维空间中"定向体积元"的自然表达——它告诉我们坐标变换是否改变定向（手性）。这正是后续学习中区分\textbf{轴矢量}（如 $\mathbf{B}$）和\textbf{极矢量}（如 $\mathbf{r}$）的数学基础。

\subsection{应避免的常见错误}

\begin{itemize}
  \item \textbf{求和指标混淆}：$\epsilon_{ijk}A_jB_k$ 中 $j,k$ 是哑指标，可以改名。但哑指标在同一项中恰好出现两次——例如 $\epsilon_{iji}A_jB_i$ 中 $i$ 出现了三次（$\epsilon$ 的第一个、第三个位置以及 $B_i$ 中），这样的写法是非法的！
  \item \textbf{$\epsilon$ 乘积次序}：$\epsilon_{ijk}\epsilon_{ilm}$ 对 $i$ 求和；$\epsilon_{jik}\epsilon_{ilm} = -\epsilon_{ijk}\epsilon_{ilm}$。
  \item \textbf{符号误判}：$(1,3,2)$ 是奇排列（一次交换），$\epsilon_{132} = -1$。
\end{itemize}

\section{三大积分定理}

\subsection{从局域到整体}

微分形式的物理定律（如 $\nabla\cdot\mathbf{E}=\rho/\varepsilon_0$）告诉我们每一点的场行为。但实验测量往往是整体的——我们测量穿过线圈的总磁通量，而不是每一点的 $\mathbf{B}$。积分定理就是连接这两种描述的桥梁。

\subsection{高斯散度定理}

\begin{theorem}[高斯定理]
对于光滑矢量场 $\mathbf{F}$ 和由光滑曲面 $\partial V$ 围成的区域 $V$，
\[
\boxed{\;\int_V (\nabla\cdot\mathbf{F})\,\dd V = \oint_{\partial V} \mathbf{F}\cdot\dd\mathbf{S}\;}.
\label{eq:gauss_div}
\]

\textbf{证明思路（一维启示）：} 一维情形下，$\int_a^b f'(x)\,\dd x = f(b)-f(a)$——内部导数的积分等于边界函数值之差。高斯定理正是这个基本定理在三维的推广。

左端是场在体积内的"源"的总强度；右端是穿过边界的"通量"。方程告诉我们：\textbf{体积内源的强度之和等于边界面流出的通量之和}。
\end{theorem}

\begin{insight}[高斯定理与电荷守恒]
在电动力学中，电流连续性方程为 $\nabla\cdot\mathbf{J} + \partial\rho/\partial t = 0$。对体积 $V$ 积分并用高斯定理：
\[
\oint_{\partial V}\mathbf{J}\cdot\dd\mathbf{S} = -\frac{\dd}{\dd t}\int_V\rho\,\dd V = -\frac{\dd Q_{\text{内}}}{\dd t}.
\]
这说的是：流出边界的电流等于内部电荷的减少率。这就是电荷守恒定律的积分形式。
\end{insight}

\begin{example}[完整算例：验证高斯定理]
取 $\mathbf{F} = (x^2, yz, z)$，$V$ 为单位立方体 $0\leq x,y,z\leq 1$。

\textbf{直接计算散度的体积分：}
\[
\nabla\cdot\mathbf{F} = 2x + z + 1.
\]
\[
\int_V \nabla\cdot\mathbf{F}\,\dd V = \int_0^1\int_0^1\int_0^1 (2x+z+1)\,\dd x\dd y\dd z = \frac52.
\]

\textbf{计算面积分：} 立方体有 6 个面。以 $x=1$ 面（法向 $\mathbf{e}_x$）为例：
\[
\int_{x=1} \mathbf{F}\cdot\dd\mathbf{S} = \int_0^1\int_0^1 (1^2)\,\dd y\dd z = 1.
\]
$x=0$ 面（法向 $-\mathbf{e}_x$）：
\[
\int_{x=0} \mathbf{F}\cdot\dd\mathbf{S} = \int_0^1\int_0^1 (0)\,\dd y\dd z = 0.
\]
类似计算其余 4 个面并求和，总面积分也为 $5/2$。高斯定理成立。$\square$
\end{example}

\subsection{斯托克斯定理}

\begin{theorem}[斯托克斯定理]
对于光滑矢量场 $\mathbf{F}$ 和以定向曲线 $\partial S$ 为边界的曲面 $S$，
\[
\boxed{\;\int_S (\nabla\times\mathbf{F})\cdot\dd\mathbf{S} = \oint_{\partial S} \mathbf{F}\cdot\dd\mathbf{l}\;}.
\label{eq:stokes}
\]

物理意义：曲面上的旋度通量等于边界上的环量。
\end{theorem}

\begin{example}[完整算例：斯托克斯定理计算环量]
$\mathbf{F} = (-y, x, 0)$，求沿 $x^2+y^2=R^2$（$z=0$）的环量。

\textbf{直接参数化：}
$\mathbf{r}(\theta) = (R\cos\theta, R\sin\theta, 0)$，$\dd\mathbf{l}=(-R\sin\theta, R\cos\theta, 0)\,\dd\theta$。
\[
\oint\mathbf{F}\cdot\dd\mathbf{l} = \int_0^{2\pi} (R^2\sin^2\theta+R^2\cos^2\theta)\,\dd\theta = 2\pi R^2.
\]

\textbf{斯托克斯定理：}
$\nabla\times\mathbf{F} = (0,0,2)$。
取圆盘 $S$（法向 $\mathbf{e}_z$）：
\[
\int_S (\nabla\times\mathbf{F})\cdot\dd\mathbf{S} = 2\cdot(\pi R^2) = 2\pi R^2.
\]
一致。注意这个结果具有几何意义：由于 $\nabla\times\mathbf{F}$ 是常数，环量正比于圆盘面积 $\pi R^2$，即正比于 $R^2$——这正是斯托克斯定理告诉我们的。
\end{example}

\subsection{格林恒等式}

由高斯定理直接导出两组恒等式，它们在 PDE 理论中极为重要。

设 $f,g$ 为光滑标量函数。

\textbf{第一格林恒等式：}
\[
\int_V (f\nabla^2 g + \nabla f\cdot\nabla g)\,\dd V = \oint_{\partial V} f\nabla g\cdot\dd\mathbf{S}.
\label{eq:green1}
\]

\textbf{第二格林恒等式：}
\[
\int_V (f\nabla^2 g - g\nabla^2 f)\,\dd V = \oint_{\partial V} (f\nabla g - g\nabla f)\cdot\dd\mathbf{S}.
\label{eq:green2}
\]

\textbf{推导：} 从高斯定理出发，令 $\mathbf{F} = f\nabla g$：
\[
\nabla\cdot(f\nabla g) = \nabla f\cdot\nabla g + f\nabla^2 g.
\]
代入高斯定理即得第一恒等式。交换 $f,g$ 后相减得第二恒等式。

\begin{insight}[格林恒等式的物理意义]
格林恒等式在电动力学中被用来证明\textbf{格林互易定理}：若 $\phi_1$ 由电荷分布 $\rho_1$ 产生，$\phi_2$ 由 $\rho_2$ 产生，则
\[
\int\rho_1\phi_2\,\dd V = \int\rho_2\phi_1\,\dd V.
\]
这个对称性关系在电容矩阵计算和电像法中非常有用。
\end{insight}

\subsection{亥姆霍兹分解}

\begin{theorem}[亥姆霍兹定理]
任何在无穷远处衰减足够快的光滑矢量场 $\mathbf{F}$ 可以唯一分解为一个无旋场和一个无散场之和：
\[
\boxed{\;\mathbf{F} = -\nabla\phi + \nabla\times\mathbf{A}\;},
\label{eq:helmholtz}
\]
其中
\[
\phi(\mathbf{r}) = \frac{1}{4\pi}\int\frac{\nabla'\cdot\mathbf{F}(\mathbf{r}')}{|\mathbf{r}-\mathbf{r}'|}\,\dd V',\qquad
\mathbf{A}(\mathbf{r}) = \frac{1}{4\pi}\int\frac{\nabla'\times\mathbf{F}(\mathbf{r}')}{|\mathbf{r}-\mathbf{r}'|}\,\dd V'.
\]

{\bfseries 这个定理为什么重要？}
\begin{itemize}
\item 在电动力学中，电场 $\mathbf{E} = -\nabla\phi - \partial\mathbf{A}/\partial t$——把 $\mathbf{E}$ 分解为库仑场（无旋）和感应场（无散在规范下）。
\item 在流体力学中，速度 $\mathbf{v} = -\nabla\phi + \nabla\times\bm{\psi}$——纵波（声波）和横波（剪切波）的分解。
\item 在地震学中，P 波（纵波，$\nabla\times\mathbf{u}=0$）和 S 波（横波，$\nabla\cdot\mathbf{u}=0$）的分离。
\end{itemize}
\end{theorem}

\section{张量初步}

\subsection{从矢量到张量：一个必要的推广}

矢量 $\mathbf{v}$ 有 3 个分量，描述力、速度等。但有些物理量需要更复杂的结构：
\begin{itemize}
  \item \textbf{转动惯量} $\mathbf{I}$：把角速度 $\bm{\omega}$ 映射为角动量 $\mathbf{L}$，$\mathbf{L} = \mathbf{I}\bm{\omega}$。如果 $\bm{\omega}$ 不沿主轴，$\mathbf{L}$ 与 $\bm{\omega}$ 方向不同。
  \item \textbf{应力} $\sigma_{ij}$：$i$ 方向的力作用在法向为 $j$ 的面上，9 个分量。
  \item \textbf{电磁场} $F^{\mu\nu}$：电场和磁场统一为四维反对称张量。
\end{itemize}

这些量的共同特征：它们给出了一个\textbf{线性映射}——把一个矢量（或更一般的张量）映射为另一个矢量（或张量）。这个映射在坐标变换下必须满足确定的变换规则，这就是张量的本质。

\subsection{张量的变换规则}

设 $V$ 是一个 $n$ 维矢量空间，$\{e_i\}$ 和 $\{e'_i\}$ 是两组基，变换矩阵为 $\Lambda^i_j = \partial x'^i/\partial x^j$。

\textbf{逆变矢量：} 分量与基矢反向变换，$V'^i = \Lambda^i_j V^j$。
\textbf{协变矢量：} 分量与基矢同向变换，$V'_i = (\Lambda^{-1})^j_i V_j$。

\textbf{$(p,q)$ 型张量：} 有 $p$ 个逆变指标和 $q$ 个协变指标，变换规则为
\[
T'^{\mu_1\cdots\mu_p}_{\nu_1\cdots\nu_q}
= \Lambda^{\mu_1}_{\alpha_1}\cdots\Lambda^{\mu_p}_{\alpha_p}
(\Lambda^{-1})^{\beta_1}_{\nu_1}\cdots(\Lambda^{-1})^{\beta_q}_{\nu_q}
\;T^{\alpha_1\cdots\alpha_p}_{\beta_1\cdots\beta_q}.
\]

\textbf{零阶张量 = 标量}，\textbf{一阶张量 = 矢量}，\textbf{二阶张量 = 矩阵（但需满足变换规则）}。

\subsection{度规张量}

曲线坐标中的弧长平方是物理不变量，可写为
\[
\dd s^2 = g_{\mu\nu}\,\dd x^\mu\dd x^\nu,
\]
其中 $g_{\mu\nu}$ 就是\textbf{度规张量}。在正交曲线坐标系中，$g_{\mu\nu} = \mathrm{diag}(h_1^2, h_2^2, h_3^2)$。

度规张量用于\textbf{升降指标}：
\[
A^\mu = g^{\mu\nu}A_\nu,\qquad A_\mu = g_{\mu\nu}A^\nu.
\]
在狭义相对论中，闵可夫斯基度规 $\eta_{\mu\nu} = \mathrm{diag}(-1,1,1,1)$ 定义了时空的几何。

\subsection{物理中最重要的二阶张量}

\textbf{转动惯量张量：}
\[
I_{ij} = \int_V \rho(\mathbf{r})(r^2\delta_{ij} - x_i x_j)\,\dd V.
\]
刚体动力学中 $\mathbf{L} = \mathbf{I}\bm{\omega}$，转动动能 $T = \frac12\bm{\omega}\cdot\mathbf{I}\bm{\omega}$。

\textbf{应力张量：}
\[
\sigma = \begin{pmatrix}
\sigma_{xx} & \sigma_{xy} & \sigma_{xz} \\
\sigma_{yx} & \sigma_{yy} & \sigma_{yz} \\
\sigma_{zx} & \sigma_{zy} & \sigma_{zz}
\end{pmatrix}.
\]
柯西应力公式 $\mathbf{T}^{(n)} = \bm{\sigma}\cdot\mathbf{n}$ 给出法向量 $\mathbf{n}$ 面上的应力矢量。

\textbf{电磁场张量：}
\[
F^{\mu\nu} = \begin{pmatrix}
0 & -E_x/c & -E_y/c & -E_z/c \\
E_x/c & 0 & -B_z & B_y \\
E_y/c & B_z & 0 & -B_x \\
E_z/c & -B_y & B_x & 0
\end{pmatrix}.
\]
麦克斯韦方程组简化为 $\partial_\mu F^{\mu\nu} = \mu_0 J^\nu$ 和 $\partial_{[\alpha}F_{\mu\nu]} = 0$。

\section{线性空间与狄拉克符号}

\subsection{从 \texorpdfstring{$\mathbb{R}^3$}{R³} 到希尔伯特空间}

你在中学就熟悉了三维矢量：可以相加、数乘、求长度和夹角。量子力学中的波函数也具有完全相同的代数结构——它们生活在\textbf{希尔伯特空间}中。

希尔伯特空间 $\mathcal{H}$ 是一个完备的内积空间。$\mathbb{R}^3$ 是希尔伯特空间的特例（有限维、实内积）；$L^2(\mathbb{R}^3)$（平方可积函数空间）是量子力学的标准舞台（无限维、复内积）。

\subsection{狄拉克的 bra-ket 符号}

狄拉克引入了一套极为优雅的符号系统：
\begin{itemize}
  \item \textbf{ket} $|\psi\rangle$：表示一个矢量（态）。
  \item \textbf{bra} $\langle\phi|$：表示一个对偶矢量（线性泛函）。
  \item \textbf{内积} $\langle\phi|\psi\rangle$：bra 作用在 ket 上得到一个复数。
  \item \textbf{外积} $|\psi\rangle\langle\phi|$：将一个矢量映射到另一个矢量的算符。
\end{itemize}

\textbf{量子力学中的物理内积：}
\[
\langle\phi|\psi\rangle = \int \phi^*(x)\psi(x)\,\dd x.
\]
模平方 $|\langle\phi|\psi\rangle|^2$ 是测量态 $|\psi\rangle$ 坍缩到 $|\phi\rangle$ 的概率。

\textbf{正交归一基：} $\langle e_i|e_j\rangle = \delta_{ij}$。
\textbf{完备性关系：}
\[
\sum_i |e_i\rangle\langle e_i| = \hat{\mathbf{1}}.
\]
\textbf{态展开：}
\[
|\psi\rangle = \sum_i |e_i\rangle\langle e_i|\psi\rangle = \sum_i c_i|e_i\rangle,
\quad c_i = \langle e_i|\psi\rangle.
\]

\begin{example}[狄拉克符号与傅里叶级数]
$\{|e_n\rangle\}$ 取为正交归一的傅里叶基 $\{e^{\mathrm{i}n\omega_0 t}/\sqrt{T}\}$。则任意周期函数可以展开为
\[
|f\rangle = \sum_n c_n |e_n\rangle,\quad
c_n = \langle e_n|f\rangle = \frac{1}{\sqrt{T}}\int_0^T f(t)e^{-\mathrm{i}n\omega_0 t}\,\dd t.
\]
这与第 3 章将学到的傅里叶级数完全一致——信号处理中的"频谱分解"在狄拉克符号中就是"希尔伯特空间中的基矢展开"。
\end{example}

\subsection{施密特正交化——完整数值算例}

给定 $\mathbb{R}^3$ 中三个矢量 $\mathbf{v}_1=(1,1,0),\;\mathbf{v}_2=(1,0,1),\;\mathbf{v}_3=(0,1,1)$。

\textbf{第 1 步：}
\[
\mathbf{e}_1 = \frac{\mathbf{v}_1}{\|\mathbf{v}_1\|} = \frac{1}{\sqrt{2}}(1,1,0).
\]

\textbf{第 2 步：}
\[
\mathbf{v}_2^\perp = \mathbf{v}_2 - (\mathbf{e}_1\cdot\mathbf{v}_2)\mathbf{e}_1 = (1,0,1) - \frac12(1,1,0) = \left(\frac12,-\frac12,1\right).
\]
\[
\mathbf{e}_2 = \frac{\mathbf{v}_2^\perp}{\|\mathbf{v}_2^\perp\|} = \frac{(1/2,-1/2,1)}{\sqrt{3/2}} = \frac{1}{\sqrt{6}}(1,-1,2).
\]

\textbf{第 3 步：}
\[
\mathbf{v}_3^\perp = \mathbf{v}_3 - (\mathbf{e}_1\cdot\mathbf{v}_3)\mathbf{e}_1 - (\mathbf{e}_2\cdot\mathbf{v}_3)\mathbf{e}_2 = \left(-\frac23,\frac23,\frac23\right).
\]
\[
\mathbf{e}_3 = \frac{\mathbf{v}_3^\perp}{\|\mathbf{v}_3^\perp\|} = \frac{1}{\sqrt{3}}(-1,1,1).
\]

验证：$\mathbf{e}_1\cdot\mathbf{e}_2 = \mathbf{e}_1\cdot\mathbf{e}_3 = \mathbf{e}_2\cdot\mathbf{e}_3 = 0$。$\square$

\section{厄米算符与谱定理}

\subsection{线性算符}

\textbf{线性算符} $\hat{A}: \mathcal{H}\to\mathcal{H}$ 满足：
\[
\hat{A}(|\psi\rangle+|\phi\rangle) = \hat{A}|\psi\rangle + \hat{A}|\phi\rangle,\qquad
\hat{A}(c|\psi\rangle) = c\hat{A}|\psi\rangle.
\]

物理中最重要的线性算符：
\begin{align*}
\hat{x}\,|\psi\rangle &= x\psi(x), &\text{位置算符}\\
\hat{p}\,|\psi\rangle &= -\mathrm{i}\hbar\frac{\dd}{\dd x}\psi(x), &\text{动量算符}\\
\hat{H}\,|\psi\rangle &= \left(-\frac{\hbar^2}{2m}\nabla^2 + V\right)\psi, &\text{哈密顿算符}
\end{align*}

\subsection{厄米算符}

\textbf{厄米算符 (Hermitian operator)} 定义为 $\hat{A}^\dagger = \hat{A}$，即
\[
\langle\phi|\hat{A}|\psi\rangle = \langle\hat{A}\phi|\psi\rangle,\qquad
\forall\,|\phi\rangle,|\psi\rangle\in\mathcal{H}.
\]

\begin{theorem}[谱定理——厄米算符的本征值问题]
若 $\hat{A}$ 是希尔伯特空间 $\mathcal{H}$ 上的厄米算符，则：
\begin{enumerate}
  \item 本征值都是\textbf{实数}。
  \item 属于不同本征值的本征矢\textbf{正交}。
  \item 本征矢构成 $\mathcal{H}$ 的\textbf{完备正交基}。
\end{enumerate}
\end{theorem}

\textbf{证明（本征值实性）：}
设 $\hat{A}|a\rangle = a|a\rangle$，则
\[
\langle a|\hat{A}|a\rangle = a\langle a|a\rangle,\qquad
\langle\hat{A}a|a\rangle = \bar{a}\langle a|a\rangle.
\]
由厄米性 $\langle a|\hat{A}|a\rangle = \langle\hat{A}a|a\rangle$，得 $a = \bar{a}$，故 $a\in\mathbb{R}$。$\square$

\begin{insight}[厄米算符 = 可观测量]
量子力学的一个基本假设：\textbf{每个可观测量对应一个厄米算符}。原因：
\begin{itemize}
  \item 本征值为实数 $\to$ 测量结果必须是实数。
  \item 本征矢正交完备 $\to$ 任意态 $|\psi\rangle$ 可以展开为 $|\psi\rangle = \sum c_n|a_n\rangle$。
  \item 测量概率 $|c_n|^2 = |\langle a_n|\psi\rangle|^2$ 满足 $\sum|c_n|^2 = 1$。
\end{itemize}
这个框架就是量子力学的"谱公理"——它告诉我们量子测量为什么是概率性的、为什么测量结果是离散的（对束缚态而言）。
\end{insight}

\subsection{完整算例：一维无限深方势阱}

\textbf{问题：} 哈密顿量 $\hat{H} = -\frac{\hbar^2}{2m}\frac{\dd^2}{\dd x^2}$，在区域 $[0,a]$ 内，边界条件 $\psi(0)=\psi(a)=0$。求本征值和本征函数。

\textbf{步骤 1：本征值方程。}
\[
-\frac{\hbar^2}{2m}\frac{\dd^2\psi}{\dd x^2} = E\psi,\quad 0<x<a.
\]

\textbf{步骤 2：通解。}
$\psi'' + k^2\psi = 0$，$k^2 = 2mE/\hbar^2$。通解 $\psi(x)=A\sin(kx)+B\cos(kx)$。

\textbf{步骤 3：边界条件。}
$\psi(0)=0 \Rightarrow B=0$。$\psi(a)=0 \Rightarrow A\sin(ka)=0 \Rightarrow ka=n\pi$，$n=1,2,3,\dots$。

\textbf{步骤 4：本征值。}
\[
k_n = \frac{n\pi}{a},\qquad
E_n = \frac{\hbar^2 k_n^2}{2m} = \frac{n^2\pi^2\hbar^2}{2ma^2}.
\]

\textbf{步骤 5：归一化本征函数。}
\[
\psi_n(x) = \sqrt{\frac{2}{a}}\sin\frac{n\pi x}{a}.
\]

\textbf{步骤 6：正交性验证。}
\[
\int_0^a \psi_n(x)\psi_m(x)\,\dd x = \frac{2}{a}\int_0^a \sin\frac{n\pi x}{a}\sin\frac{m\pi x}{a}\,\dd x = \delta_{nm}.
\]

\textbf{物理解读：}
\begin{itemize}
  \item 基态 $E_1 = \pi^2\hbar^2/(2ma^2)$——零点能，不确定性原理的直接后果。
  \item 激发态 $E_n\sim n^2$——能级间隔随 $n$ 增大。
  \item 波函数节点数 $=n-1$——节点越多，能量越高。
\end{itemize}

\subsection{对易关系与不确定原理}

\textbf{对易子：} $[\hat{A},\hat{B}] = \hat{A}\hat{B} - \hat{B}\hat{A}$。如果 $[\hat{A},\hat{B}]=0$，称 $\hat{A}$ 和 $\hat{B}$ 对易，它们有共同的本征态。

\textbf{基本对易关系：}
\[
[\hat{x},\hat{p}] = \mathrm{i}\hbar\hat{\mathbf{1}}.
\]

\textbf{不确定原理：}
\[
\Delta A\,\Delta B \geq \frac12|\langle[\hat{A},\hat{B}]\rangle|.
\]
代入 $[\hat{x},\hat{p}]=\mathrm{i}\hbar$：
\[
\Delta x\,\Delta p \geq \frac{\hbar}{2}.
\]

\begin{insight}[不确定原理的几何解释]
不确定原理不是仪器的缺陷，而是量子态几何性质的体现。两个算符不对易意味着它们的本征态“指向不同方向”——你不能同时精确确定它们。$\Delta x\,\Delta p \geq \hbar/2$ 的数学根源在傅里叶分析中：时宽-带宽积受 $\Delta t\,\Delta\omega \geq 1/2$ 约束（见第 3 章）。位置和动量正是通过傅里叶变换联系的。
\end{insight}

\section{应用专题：电磁场与量子力学中的数学结构}

\subsection{电磁场中高斯定理的应用}

回顾 \S1.6 的高斯定理。在电动力学中，我们有：
\[
\oint_{\partial V}\mathbf{E}\cdot\dd\mathbf{S} = \frac{Q_{\text{内}}}{\varepsilon_0}.
\]

这个方程本质上不是"物理定律"而是库仑定律的几何等价形式。利用它，可以绕过复杂的积分直接求出对称电荷分布的电场：

\begin{itemize}
  \item 球对称电荷分布：$E(r) = Q_{\text{内}}(r)/(4\pi\varepsilon_0 r^2)$。
  \item 无限长线电荷：$E(\rho) = \lambda/(2\pi\varepsilon_0\rho)$。
  \item 无限大平板：$E = \sigma/(2\varepsilon_0)$。
\end{itemize}

这些结果只需要计算 $Q_{\text{内}}$ 和面积分——不需要解任何微分方程！

\subsection{从谱定理到量子测量}

量子力学的基本假设可以用厄米算符的语言统一表述：

\begin{enumerate}
  \item 系统的态由希尔伯特空间中的矢量 $|\psi\rangle$ 描述（$\langle\psi|\psi\rangle=1$）。
  \item 可观测物理量由厄米算符 $\hat{A}$ 表示。
  \item 测量 $\hat{A}$ 得到本征值 $a_n$，概率为 $|\langle a_n|\psi\rangle|^2$。
  \item 测量后态坍缩到 $|a_n\rangle$。
  \item 态的时间演化由薛定谔方程 $\mathrm{i}\hbar\partial_t|\psi\rangle = \hat{H}|\psi\rangle$ 决定。
\end{enumerate}

这些假设看似抽象，但每一步都有其数学必然性——第 1-4 步由谱定理保证，第 5 步由 $\hat{H}$ 的厄米性保证概率守恒。

\section{本章公式速查}

\begin{table}[h]
\centering
\caption{第 1 章核心公式}
\begin{tabular}{c|c}
\toprule
概念 & 公式 \\
\midrule
拉梅系数 & $h_i = |\partial\mathbf{r}/\partial q_i|$ \\
弧长元 & $\dd s^2 = h_1^2\dd q_1^2 + h_2^2\dd q_2^2 + h_3^2\dd q_3^2$ \\
梯度 & $\nabla f = \sum (1/h_i)(\partial f/\partial q_i)\mathbf{e}_i$ \\
散度 & $\nabla\cdot\mathbf{A} = (1/h_1h_2h_3)\sum\partial_i(h_1h_2h_3A_i/h_i)$ \\
旋度 & $\nabla\times\mathbf{A} = \sum (1/h_jh_k)[\partial_j(h_kA_k)-\partial_k(h_jA_j)]\mathbf{e}_i$ \\
拉普拉斯 & $\nabla^2 f = (1/h_1h_2h_3)\sum\partial_i((h_1h_2h_3/h_i^2)\partial_i f)$ \\
$\epsilon$-$\delta$ 缩并 & $\epsilon_{ijk}\epsilon_{ilm} = \delta_{jl}\delta_{km} - \delta_{jm}\delta_{kl}$ \\
高斯定理 & $\int_V\nabla\cdot\mathbf{F}\,\dd V = \oint_{\partial V}\mathbf{F}\cdot\dd\mathbf{S}$ \\
斯托克斯定理 & $\int_S(\nabla\times\mathbf{F})\cdot\dd\mathbf{S} = \oint_{\partial S}\mathbf{F}\cdot\dd\mathbf{l}$ \\
亥姆霍兹分解 & $\mathbf{F} = -\nabla\phi + \nabla\times\mathbf{A}$ \\
第一格林恒等式 & $\int_V(f\nabla^2g + \nabla f\cdot\nabla g) = \oint f\nabla g\cdot\dd\mathbf{S}$ \\
厄米算符 & $\langle\phi|\hat{A}|\psi\rangle = \langle\hat{A}\phi|\psi\rangle$ \\
不确定原理 & $\Delta A\,\Delta B \geq \frac12|\langle[\hat{A},\hat{B}]\rangle|$ \\
\bottomrule
\end{tabular}
\end{table}

\section{本章概念图}

\begin{center}
\fbox{\parbox{0.88\textwidth}{
\textbf{物理问题} $\xrightarrow{\text{选择坐标系}}$ \textbf{拉梅系数} $\xrightarrow{\text{微分算符}}$ \textbf{梯度/散度/旋度} \\
$\downarrow$ \\
\textbf{积分定理}（高斯、斯托克斯）$\xrightarrow{\text{整体-局部}}$ \textbf{亥姆霍兹分解} \\
$\downarrow$ \\
\textbf{张量} $\xrightarrow{\text{坐标变换下不变}}$ \textbf{线性空间 + 厄米算符} \\
$\downarrow$ \\
\textbf{谱定理} $\xrightarrow{\text{本征值}}$ \textbf{可观测量与量子测量}
}}
\end{center}

\section*{习题}

\textbf{基础题（巩固概念）——建议用时 45 分钟}

\begin{enumerate}
  \item 写出柱坐标系 $(\rho,\phi,z)$ 中梯度、散度、旋度和拉普拉斯算子的显式表达式。然后计算 $\nabla\cdot(\rho\mathbf{e}_\rho)$。
  \item 用指标法证明 $\nabla\cdot(\nabla\times\mathbf{A}) = 0$ 和 $\nabla\times(\nabla f) = 0$。
  \item 计算三重积分：直接验证 $\displaystyle\int_0^1\int_0^1\int_0^1 (2x+y+z)\,\dd x\dd y\dd z$。然后验证高斯定理对 $\mathbf{F}=(2x,y,z)$ 在单位立方体上成立。
  \item 判断下列矢量是否线性无关：\\ (1) $(1,2,3),(4,5,6),(7,8,9)$；(2) $(1,0,1),(0,1,1),(1,1,0)$。
  \item 求 $f(z)=z^2$ 在 $z=1+\mathrm{i}$ 处的导数（回顾第 2 章内容，此处验证第 1 章线性空间概念的普遍性）。
\end{enumerate}

\textbf{提高题（应用分析）——建议用时 60 分钟}

\begin{enumerate}
  \setcounter{enumi}{5}
  \item 求球坐标中 $\mathbf{A}=r^n\mathbf{e}_r$ 的散度。$n$ 取何值时散度为零？解释这一结果对静电场中点电荷场的意义。
  \item 设 $\mathbf{F}=(y,-x,z)$，用斯托克斯定理求沿 $x^2+y^2=1$（$z=0$）的环量。
  \item 证明 $\mathbf{A}\times(\mathbf{B}\times\mathbf{C}) + \mathbf{B}\times(\mathbf{C}\times\mathbf{A}) + \mathbf{C}\times(\mathbf{A}\times\mathbf{B}) = 0$（雅可比恒等式，用 $\epsilon_{ijk}$ 证明）。
  \item 证明转动惯量张量 $I_{ij}$ 是对称的。对长方体 $0\leq x\leq a,0\leq y\leq b,0\leq z\leq c$，写出 $I_{ij}$ 的显式表达式。
  \item 一维谐振子的哈密顿量 $\hat{H} = -\frac{\hbar^2}{2m}\frac{\dd^2}{\dd x^2} + \frac12 m\omega^2 x^2$。\\ 证明升降算符 $\hat{a} = \sqrt{m\omega/(2\hbar)}(\hat{x} + \mathrm{i}\hat{p}/(m\omega))$ 满足 $[\hat{a},\hat{a}^\dagger]=1$，且本征值为 $E_n = (n+1/2)\hbar\omega$。
\end{enumerate}

\textbf{挑战题（综合拓展）——建议用时 90 分钟}

\begin{enumerate}
  \setcounter{enumi}{10}
  \item 证明角动量算符 $\hat{L}_i = \epsilon_{ijk}\hat{x}_j\hat{p}_k$ 满足 $[\hat{L}_i,\hat{L}_j] = \mathrm{i}\hbar\epsilon_{ijk}\hat{L}_k$（需用到 $[\hat{x}_i,\hat{p}_j]=\mathrm{i}\hbar\delta_{ij}$ 和 $\epsilon$-$\delta$ 公式）。
  \item 写出四维势 $A^\mu = (\phi/c,\mathbf{A})$，构造 $F_{\mu\nu} = \partial_\mu A_\nu - \partial_\nu A_\mu$。验证 $\partial_\mu F^{\mu\nu} = \mu_0 J^\nu$ 给出麦克斯韦方程组中的两个含源方程。
  \item 用亥姆霍兹分解证明：无界区域内若 $\nabla\cdot\mathbf{F}=0$ 且 $\nabla\times\mathbf{F}=0$，则 $\mathbf{F}=\mathbf{0}$。讨论这一结论在电磁腔本征模分析中的意义。
  \item 设 $\hat{A}$ 和 $\hat{B}$ 都是厄米算符。证明 $\mathrm{i}[\hat{A},\hat{B}]$ 也是厄米算符。这为什么能保证不确定原理 $\Delta A\,\Delta B \geq \frac12|\langle[\hat{A},\hat{B}]\rangle|$ 的右端是实数？
  \item 编程：在 $\mathbb{R}^3$ 中取 4 个不完全线性无关的矢量，用 Gram-Schmidt 正交化并可视化。验证数值误差对正交性的影响。
  \item （研究导向）阅读资料，了解"轴矢量"和"极矢量"在宇称变换下的区别。解释为什么 $\mathbf{B}$ 是轴矢量而 $\mathbf{E}$ 是极矢量。在 $\epsilon_{ijk}$ 的框架下，轴矢量的变换行为是怎样的？
\end{enumerate}
